\documentclass[12pt, a4paper, oneside]{ctexart}
\usepackage{amsmath, amsthm, amssymb, bm, color, framed, graphicx, hyperref, mathrsfs}
\usepackage{booktabs}
\usepackage{geometry}

\title{\textbf{4月27日作业}}
\author{韩岳成 524531910029}
\date{\today}
\linespread{1.5}
\definecolor{shadecolor}{RGB}{241, 241, 255}
\newcounter{problemname}
\newenvironment{problem}{\begin{shaded}\stepcounter{problemname}\par\noindent\textbf{题目\arabic{problemname}. }}{\end{shaded}\par}
\newenvironment{solution}{\par\noindent\textbf{解答. }}{\par}
\newenvironment{note}{\par\noindent\textbf{题目\arabic{problemname}的注记. }}{\par}
\everymath{\displaystyle}

\begin{document}

\maketitle

\begin{problem}
        设 \(G\) 是简单图，证明：
\[
\frac{1+\operatorname{diam}G}{3}\leqslant\gamma(G)\leqslant\nu-\left\lceil\frac{\operatorname{diam}G}{2}\right\rceil
\]
\end{problem}

\begin{solution}


下界：取一条长度为直径 \(d = \operatorname{diam} G\) 的最短路径 \(P: v_0, v_1, \dots, v_d\)。设 \(D\) 是 \(G\) 的一个最小支配集，\(|D| = \gamma(G)\)。对每个 \(i \in \{0,1,\dots,d\}\)，因 \(D\) 支配所有顶点，存在顶点 \(w_i \in D\) 使得 \(w_i\) 与 \(v_i\) 相邻或 \(w_i = v_i\)。考虑集合 \(S = \{w_i : i = 0,\dots,d\}\)，这是 \(D\) 的一个子集（可能有重复）。我们断言：若 \(w_i = w_j = w\)，则 \(|i - j| \le 2\)。否则，若 \(|i - j| \ge 3\)，则从 \(v_0\) 到 \(v_d\) 可走路径
\[
v_0 \to \cdots \to v_i \to w \to v_j \to \cdots \to v_d,
\]
其长度为 \(i + 1 + 1 + (d - j) = d - (j - i) + 2 \le d - 1\)，与 \(d\) 是直径矛盾。因此，每个 \(w \in D\) 至多支配 \(P\) 上的三个顶点（因为索引差不超过 2 意味着这些顶点包含在某个长度为 2 的区间内，最多三个点）。于是 \(P\) 上的 \(d+1\) 个顶点被 \(D\) 中的元素覆盖，每个元素覆盖不超过 3 个，故
\[
\gamma(G) = |D| \ge \left\lceil \frac{d+1}{3} \right\rceil \ge \frac{\operatorname{diam}G+1}{3}.
\]
下界得证。

上界：仍取直径路径 \(P: v_0, v_1, \dots, v_d\)。由于 \(P\) 是 \(G\) 中所有最短路径中最长的一条，\(P\) 上任意两个非相邻顶点之间无边（否则会缩短 \(v_0\) 与 \(v_d\) 的距离）。取 \(P\) 上所有下标为偶数的顶点，构成集合 \(I = \{v_0, v_2, v_4, \dots\}\)，则 \(I\) 中任意两点在 \(P\) 上距离至少为 2，且因非相邻顶点之间无边，它们在 \(G\) 中也不相邻，故 \(I\) 是独立集。\(I\) 的大小为 \(\lceil (d+1)/2 \rceil\)。于是其补集 \(V(G) \setminus I\) 是支配集（因为 \(I\) 中每个顶点仅与 \(V(G) \setminus I\) 中的顶点相邻，且 \(V(G) \setminus I\) 自身显然被支配），从而
\[
\gamma(G) \le \nu - |I| = \nu - \left\lceil \frac{d+1}{2} \right\rceil \le \nu - \left\lceil \frac{\operatorname{diam}G}{2} \right\rceil,
\]


综上，原不等式成立。∎
\end{solution}

\begin{problem}
证明：

1.极小点覆盖集的补集是极大点独立集；

2. 极大点独立集的补集是极小点覆盖集；

3.对于任意一个图$G:\alpha(G)+\beta(G)=\nu(G)$。
\end{problem}

\begin{solution}

1.
设 \(C\) 是极小点覆盖，令 \(I = V \setminus C\)。  

独立性：若存在 \(u,v \in I\) 且 \(uv \in E\)，则边 \(uv\) 的两个端点都不在 \(C\) 中，与 \(C\) 是点覆盖矛盾。故 \(I\) 是点独立集。  

极大性：任取 \(v \in C\)，考虑 \(I \cup \{v\}\)。因 \(C\) 极小，\(C \setminus \{v\}\) 不是点覆盖，故存在边 \(e = xy\) 不被 \(C \setminus \{v\}\) 覆盖，但被 \(C\) 覆盖。由于只有 \(v\) 被移除，该边必以 \(v\) 为一端点，另一端点 \(u \notin C\)（否则边仍被覆盖）。于是 \(u \in I\)，且 \(v\) 与 \(u\) 相邻，故 \(I \cup \{v\}\) 含相邻点，不是独立集。因此 \(I\) 是极大点独立集。

2. 
设 \(I\) 是极大独立集，令 \(C = V \setminus I\)。  

覆盖性：任取边 \(uv\)，若 \(u,v \notin C\)，则 \(u,v \in I\)，与 \(I\) 独立矛盾。故至少一端在 \(C\) 中，\(C\) 是点覆盖集。  

极小性：任取 \(v \in C\)，考虑 \(C \setminus \{v\}\)。因 \(I\) 极大，\(I \cup \{v\}\) 不是独立集，故存在 \(u \in I\) 使 \(uv \in E\)。此时 \(u \notin C \setminus \{v\}\)（因为 \(u \in I\)），且 \(v\) 被移除，故边 \(uv\) 两端均不在 \(C \setminus \{v\}\) 中，即 \(C \setminus \{v\}\) 不覆盖 \(uv\)，不是点覆盖。因此 \(C\) 是极小点覆盖。



3.   

若 \(S\) 是最大点独立集，则 \(|S| = \alpha(G)\)，且 \(V \setminus S\) 一定是最小点覆盖集。于是  
\[
|V \setminus S| = \beta(G),
\]
从而  
\[
\alpha(G) + \beta(G) = |S| + |V \setminus S| = |V| = \nu(G).
\]


\end{solution}

\begin{problem}
证明：任一个二部图$G$有一个至少含$\frac{\varepsilon(G)}{\Delta(G)}$条边的匹配.并由此证明：若$H$是完全二部图$K_{n,n}$的一个子图，且$\varepsilon(H)>(k-1)n$,则$H$中存在至少含有$k$条边的匹配.
\end{problem}

\begin{solution}

(1) 设二部图 \(G\) 的边数为 \(\varepsilon(G)\)，最大度为 \(\Delta(G)\)。令 \(M\) 为 \(G\) 的一个最大匹配，并记 \(\alpha(G) = |M|\)。由于在二部图中，$\alpha(G) = \beta'(G)$，故 \(G\) 存在一个最小点覆盖 \(C\) 满足 \(|C| = \alpha(G)\)。点覆盖 \(C\) 包含了 \(G\) 中所有边的至少一个端点，因此
\[
\varepsilon(G) \le \sum_{v\in C} \deg(v) \le |C| \cdot \Delta(G) = \alpha(G) \cdot \Delta(G).
\]
从而得到 \(\alpha(G) \ge \frac{\varepsilon(G)}{\Delta(G)}\)。由于匹配边数 \(\alpha(G)\) 是整数，这说明 \(G\) 中存在至少 \(\left\lceil \frac{\varepsilon(G)}{\Delta(G)} \right\rceil\) 条边的匹配，当然也就存在至少 \(\frac{\varepsilon(G)}{\Delta(G)}\) 条边的匹配。

(2) 设 \(H\) 是完全二部图 \(K_{n,n}\) 的子图，且 \(\varepsilon(H) > (k-1)n\)。在 \(K_{n,n}\) 中，每个顶点的度数不超过 \(n\)，因此其子图 \(H\) 的最大度满足 \(\Delta(H) \le n\)。由 (1) 的结论，\(H\) 中存在匹配 \(M\) 满足
\[
|M| \ge \frac{\varepsilon(H)}{\Delta(H)} > \frac{(k-1)n}{n} = k-1.
\]
因为 \(|M|\) 是整数，\(|M| > k-1\) 意味着 \(|M| \ge k\)。故 \(H\) 中存在至少含有 \(k\) 条边的匹配。
\end{solution}

\begin{problem}
证明:图$G$是二部图当且仅当对$G$的每个满足$\delta(H)>0$的子图$H$均有$\alpha(H)=\beta^{\prime}(H)$
\end{problem}

\begin{solution}
必要性:

设 \(G\) 是二部图，\(H\) 是 \(G\) 的任意一个满足 \(\delta(H) > 0\) 的子图（即无孤立顶点）。因为二部图的子图仍是二部图，且 \(H\) 无孤立点，根据已知结论直接得到 \(\alpha(H) = \beta'(H)\)。必要性得证。

充分性: 

假设对 \(G\) 的每个满足 \(\delta(H) > 0\) 的子图 \(H\)，均有 \(\alpha(H) = \beta'(H)\)。  

如果 \(G\) 不是二部图，则 \(G\) 必含有一个奇圈，记为 \(C = C_{2k+1} \ (k \ge 1)\)。取子图 \(H = C\)，显然 \(\delta(C) = 2 > 0\)。

对于这个奇圈：最大独立集大小 \(\alpha(C) = k\)，最小边覆盖大小 \(\beta'(C) = k+1\)。

因此 \(\alpha(C) \neq \beta'(C)\)，与假设矛盾。所以 \(G\) 不含奇圈，即 \(G\) 是二部图。

综上，命题得证。
\end{solution}

\begin{problem}
请将本章学习的6个量在一般图和二部图中的关系汇总成一张表，并总结每组关系的应用。
\end{problem}
\begin{solution}


\textbf{符号说明：}
\begin{itemize}
    \item $\nu(G)$：顶点数
    \item $\alpha(G)$：点独立数（最大独立集大小）
    \item $\beta(G)$：点覆盖数（最小点覆盖大小）
    \item $\gamma(G)$：点支配数（最小支配集大小）
    \item $\alpha'(G)$：边独立数（最大匹配大小）
    \item $\beta'(G)$：边覆盖数（最小边覆盖大小）
    \item $\gamma'(G)$：边支配数（最小边支配集大小）
\end{itemize}


\begin{table}[h]
\centering
\begin{tabular}{@{}lcc@{}}
\toprule
\textbf{参数关系} & \textbf{一般图} & \textbf{二部图} \\
\midrule
$\alpha + \beta$ & $= \nu$ & $= \nu$ \\
$\alpha' + \beta'$ & $= \nu$（无孤立点） & $= \nu$\\
$\alpha'$ 与 $\beta$ & $\alpha' \le \beta$ & $\alpha' = \beta$\\
$\alpha$ 与 $\beta'$ & $\alpha \le \beta'$ & $\alpha = \beta'$ \\
$\gamma$ 与 $\alpha$、$\beta$ & 无简单等式，但有界 & 同一般图 \\
$\gamma'$ 与 $\alpha'$ & $\alpha' \le \gamma' \le 2\alpha'$（无孤立点） & 同一般图 \\
$\gamma'$ 与 $\beta'$ & 无固定大小关系，例：星图 $\gamma'=1,\ \beta'=n$ & 同一般图 \\
\bottomrule
\end{tabular}
\end{table}

\subsection*{每组关系的应用总结}
\begin{enumerate}
    \item \textbf{点独立与点覆盖}：$\alpha + \beta = \nu$ 提供了相互转换的基础。应用：顶点覆盖问题的对偶是独立集问题；在参数算法中，若一个参数小，另一个必大；用于近似算法设计（如贪心覆盖与极大独立集的2倍近似比）。

    \item \textbf{边独立与边覆盖}：$\alpha' + \beta' = \nu$（无孤立点）表明最大匹配与最小边覆盖互补。应用：求解最小边覆盖可先求最大匹配，再加未覆盖点的关联边；在二分图匹配中用于指派问题、工作分配。

    \item \textbf{最大匹配与最小点覆盖}：仅适用于二部图，$\alpha' = \beta$。应用：这是组合优化中最重要的定理之一，将匹配问题转化为点覆盖，进而转化为网络流；用于计算二分图的最小顶点覆盖、最大独立集（$\alpha = \nu - \beta$）；在矩阵中求最小行/列覆盖、在图像分割、推荐系统中有广泛应用。

    \item \textbf{点支配集}：虽无简单对偶，但下界 $\gamma \ge \frac{\nu}{1+\Delta}$（$\Delta$ 为最大度）等用于无线传感器网络选址、设施选址问题；经典贪心算法给出 $O(\log \nu)$ 近似比。

    \item \textbf{边支配集与匹配}：$\alpha' \le \gamma' \le 2\alpha'$ 说明最小边支配集大小介于最大匹配和其两倍之间。应用：用于网络边监控、无线网络中边覆盖的节能问题；可通过最大匹配构造边支配集（取匹配边及其邻边），近似比为2。

\end{enumerate}


\end{solution}
\end{document}